% Ch. 13 : configuration, state, réalité
\documentclass[border=6pt]{standalone}
\usepackage{coursfig}
\begin{document}
\begin{tikzpicture}[x=1mm,y=1mm]
  \tikzset{src/.style={box=#1, text width=46mm, minimum height=20mm}}
  \node[src=Enc]  (c) at (-60,0) {\ico[Enc]{file-code}\enspace\textbf{Configuration}\sub{\cmd{*.tf} : l'état désiré, dans Git}};
  \node[src=Ocre] (s) at (0,0)   {\ico[Ocre]{database}\enspace\textbf{State}\sub{\cmd{terraform.tfstate} : ce que}\sub{Terraform croit avoir créé}};
  \node[src=Sea]  (r) at (60,0)  {\ico[Sea]{cloud}\enspace\textbf{Réalité}\sub{ce que l'API rapporte}\sub{vraiment}};
  \node[keybox=Ink, text width=40mm, minimum height=13mm, font=\ttfamily\large\bfseries] (p) at (0,-38) {terraform plan};
  \node[keybox=Brick, text width=40mm, minimum height=13mm, font=\ttfamily\large\bfseries] (a) at (0,-70) {terraform apply};

  \draw[flow=Enc] (c.south) |- node[elabel, pos=.25, left=1mm, fill=none, align=right]{différence ?} (p.west);
  \draw[flow=Ocre] (s) -- (p);
  \draw[flow=Sea] (r.south) |- node[elabel, pos=.25, right=1mm, fill=none, align=left]{refresh : le state est\\confronté à la réalité} (p.east);
  \draw[flow] (p) -- node[elabel, right=2mm, fill=none]{plan validé} (a);
  \node[note, anchor=west, text width=60mm] at ([xshift=6mm]a.east) {exécute les changements via l'API,\\puis met à jour le state};
\end{tikzpicture}
\end{document}
